\section{Closed-Walk Moment Obstructions for General Period}\label{sec:general-period}

The local square identity in Section~\ref{sec:eight-barrier} does not depend on the period. We now
derive the first three moments for every legal periodic phase and prove
Theorem~\ref{thm:general-period-moment-obstruction}.

\subsection{From closed walks to flux monomials}

The possible steps of \eqref{eq:periodic-operator} are \(-2,-1,+1,+2\). A step of length one has
weight one. A \(+2\) step from \(j\) has weight \(\tau_j\), and a \(-2\) step from
\(j\) has weight \(\tau_{j-2}\).

Consider a closed step word and cancel repeated \(\tau\) factors modulo two.
The remaining indices occur in even number; write them as

\begin{equation}
r_{1}<r_{2}<\ldots<r_{2s}.
\end{equation}

Using \(Q_h=\tau_h \tau_{h+1}\) and telescoping,

\begin{equation}
\label{eq:tau-product-from-flux}
\begin{aligned}
&\prod_j \tau_{r_j} \\
&=\prod_{\ell=1}^{s}\prod_{h=r_{2\ell-1}}^{r_{2\ell}-1}Q_h.
\end{aligned}
\end{equation}

Thus every signed closed walk becomes a monomial in interval products of \(Q\).
The identity is on the integer lattice, before reduction modulo a period, so
it also handles short-cell collisions once cyclic indices are imposed.

\subsection{Exact closed-word collection}

For completeness, the collection has the following finite dynamic-programming
recurrence on the integer lattice. For a starting residue \(r\), set

\begin{equation}
\label{eq:walk-recurrence}
\begin{aligned}W_0^{(r)}(j)&=\begin{cases}1,&j=r,\\0,&j\ne r,\end{cases}\\
W_{\ell+1}^{(r)}(j)&=\sum_{\delta\in\{-2,-1,1,2\}}W_{\ell}^{(r)}(j-\delta)w(j-\delta,\delta).
\end{aligned}
\end{equation}

where

\begin{equation}
w(i,-1)=w(i,+1)=1,   w(i,+2)=\tau_i,   w(i,-2)=\tau_{i-2}.
\end{equation}

After every multiplication, pairs of equal \(\tau\) factors are cancelled and
the remaining factors are converted to \(Q\) interval products by \eqref{eq:tau-product-from-flux}. Then

\begin{equation}
\label{eq:walk-moment}
M_k=\sum_{r=0}^{p-1} W_{2k}^{r}(r).
\end{equation}

Thus every coefficient below is obtained by additions and multiplications in
the integral group algebra of the free sign variables; no floating-point
spectral calculation enters the collection.

There are respectively 4, 36, and 430 closed step words of lengths two, four,
and six. Collecting their flux monomials up to translation gives

\begin{table}[tbp]
\centering
\caption{Closed-walk contributions from one starting residue.}
\label{tab:closed-walk-contributions}
\begingroup
\renewcommand{\arraystretch}{1.12}
\begin{tabular}{ll}
\toprule
length & contribution from one starting residue \\
\midrule
2 & \(4\) \\
4 & \(28+8Q_i\) \\
6 & \(238+156Q_i+24Q_iQ_{i+1}+12Q_iQ_{i+2}\) \\
\bottomrule
\end{tabular}
\endgroup
\end{table}

We briefly justify the collection. At length two, each of the four steps is
followed by its reverse. At length four, the 36 zero-sum words reduce via
\eqref{eq:tau-product-from-flux} to 28 constant terms and eight translates of a single \(Q_i\). At length
six, applying \eqref{eq:tau-product-from-flux} to the 430 zero-sum words and grouping equal translated
interval products leaves only the four displayed monomial types. This is a
finite symbolic identity in the free \(Q\) variables, not a period-by-period
numerical observation.

Summing over the \(p\) starting residues yields

\begin{equation}
\label{eq:general-flux-moments}
\begin{aligned}
M_{1}&=4p, \\
M_{2}&=28p+8 \sum_i Q_i, \\
M_{3}&=238p+156 \sum_i Q_i \\
&\quad +24 \sum_i Q_iQ_{i+1}+12 \sum_i Q_iQ_{i+2}.
\end{aligned}
\end{equation}

Let \(I_i=(1+Q_i)/2\). Then

\begin{equation}
d=\sum_i I_i,       a=\sum_i I_iI_{i+1},       b=\sum_i I_iI_{i+2}.
\end{equation}

Substitute \(Q_i=2I_i-1\) in \eqref{eq:general-flux-moments} and collect. The result is

\begin{equation}
\label{eq:general-defect-moments}
\begin{aligned}
M_{1}&=4p, \\
M_{2}&=20p+16d, \\
M_{3}&=118p+168d+96a+48b.
\end{aligned}
\end{equation}

Because all sums are cyclic, \eqref{eq:general-defect-moments} remains valid for \(p=1,2,3,4\), even when
the offsets one and two coincide modulo \(p\); multiplicities are retained by
the Laurent-matrix construction.

\subsection{Necessary conditions at the eight-barrier}

From \eqref{eq:general-defect-moments},

\begin{equation}
\label{eq:general-excesses}
\begin{aligned}
F_{1}&=M_{2}-8M_{1}=16d-12p, \\
F_{2}&=M_{3}-8M_{2}=-42p+40d+96a+48b.
\end{aligned}
\end{equation}

If \(R(Q)\leq 8\), Lemma~\ref{lem:moment-barrier} requires \(F_{1}\leq 0\) and \(F_{2}\leq 0\). Therefore

\begin{equation}
\label{eq:general-necessary-conditions}
\begin{aligned}
d&\leq \frac{3p}{4}, \\
40d+96a+48b&\leq 42p.
\end{aligned}
\end{equation}

This proves Theorem~\ref{thm:general-period-moment-obstruction}.

The \(4\), \(36\), and \(430\) word collections are a finite computer-assisted
symbolic calculation. Recurrences \eqref{eq:walk-recurrence}-\eqref{eq:walk-moment} specify the complete proof
algorithm, and the exact checker named in Appendix~\ref{app:computational-protocol} reproduces the grouped
table, formulas \eqref{eq:general-flux-moments}-\eqref{eq:general-excesses}, and the period-eight values
\(F_{4}=5504\), \(F_{6}=64336\), and \(F_{9}=2872096\) used in Theorem~\ref{thm:eight-barrier-trichotomy}.

The first inequality limits defect density. The second also penalizes local
clustering, assigning greater weight to adjacent pairs than to distance-two
pairs. The converse is false in general: satisfying \eqref{eq:general-necessary-conditions}, or observing
finitely many nonpositive \(F_k\), does not establish \(R(Q)\leq 8\).

\subsection{Relation to the period-eight mechanism}

At \(p=8\), equations \eqref{eq:general-defect-moments}-\eqref{eq:general-excesses} specialize to \eqref{eq:period-eight-first-moment}-\eqref{eq:period-eight-second-excess}. For high defect
density the first two excesses already cross the barrier. Low-density phases
can evade these short tests, and the two-defect separation hierarchy in
Section~\ref{sec:eight-barrier} shows that progressively longer closed walks are then required.
This motivates the adaptive hierarchy used in the bounded low-period theorem.
